\section{Related work}
\label{sec:related}

\subsection{Document management and content-based organisation}

Enterprise content management addresses the governance of documents over
time, including how they are described, stored and shared
\citep{paivarinta-ecm-2005,tyrvaainen-ecm-2006}. Metadata quality and tagging
affect what can be found \citep{park-metadata-quality-2009,golder-tagging-2006};
SKOS and Dublin Core supply conventions for vocabularies and descriptions
\citep{skos,dcmi-terms}. Text categorisation addresses the assignment of
documents to classes \citep{sebastiani2002machine}, while layout-aware parsing
recovers text and structure for downstream operations
\citep{zhong-publaynet-layout-2019,huang-layoutlmv3-docai-2022,mineru}.
In HMRS, placement is not only an ingest outcome: it determines which
collection and shelf show the document to a later query agent. The workflow
uses recovered content to assist placement and description, and checks
whether managed files remain reachable after selected lifecycle operations.
These checks concern whether a stored document is visible to search after
a write. SPFresh addresses a complementary problem: maintaining a vector
index incrementally
\citep{xu2023spfresh}.

\subsection{Collection selection and structured retrieval}

Distributed information retrieval distinguishes collection selection from
ranking and merging results \citep{callan1995searching,si2003redde,shokouhi-federated-2011}.
CORI and ReDDE use representations of collection contents to decide where
to search; HMRS likewise represents collections, but exposes their managed
descriptions and membership to an agent that can then inspect descriptions
of documents on the selected shelf. A known scope can instead be passed to
pinned search. Thus explicit scoping is not an alternative to BM25, dense
retrieval or rank fusion: those methods can rank candidates once the search
scope has been selected \citep{bm25,dpr,rrf-cormack,bge-m3}.

\begin{table*}[!t]
\centering
\small
\setlength{\tabcolsep}{3.5pt}
\renewcommand{\arraystretch}{1.02}
\caption{Representative systems compared by the structure they search and
how that structure is maintained and used. A dash indicates that the table does not
characterise the system on that dimension.}
\label{tab:related-comparison}
\begin{tabular}{@{}>{\raggedright\arraybackslash}p{0.13\textwidth}>{\raggedright\arraybackslash}p{0.245\textwidth}>{\raggedright\arraybackslash}p{0.21\textwidth}>{\raggedright\arraybackslash}p{0.215\textwidth}>{\raggedright\arraybackslash}p{0.13\textwidth}@{}}
\toprule
System & Hierarchy and scope & Construction and lifecycle & Evidence handling & Access surface\\
\midrule
RAPTOR \citep{raptor} & cluster tree over document text & built offline at index time & --- & tree retrieval\\
GraphRAG \citep{graphrag} & entity--relation graph, community summaries & pre-generated offline & --- & graph retrieval\\
LightRAG \citep{lightrag} & low- and high-level graph retrieval & incremental graph updates & --- & graph retrieval\\
DeepRead \citep{deepread} & per-document structured Markdown, section coordinates & structure built at ingestion & agent locates then reads sections & agentic search\\
EPRAG \citep{kinalioglu-kbs-eprag-2026} & multi-source enterprise RAG & --- & --- & enterprise retrieval\\
HMRS (this work) & managed collection--document--segment hierarchy & content-informed placement and lifecycle checks & source reading; optional passage gate & self-hosted; MCP layer\\
\bottomrule
\end{tabular}
\end{table*}

Hierarchies and graphs need not represent the same objects. RAPTOR clusters
text into a tree for retrieval \citep{raptor}, and GraphRAG and LightRAG
provide graph-based access to content \citep{graphrag,lightrag}. RMR-SST
represents engineering documents as semantic trees
\citep{li-kbs-rmrsst-2026}. DeepRead allows agents to locate and read
sections using document structure \citep{deepread}. HMRS shares the goal of
structure-aware access, but its upper levels are managed collections and
the documents assigned to them. It uses that maintained organisation to
choose where to look and which documents to open before reaching particular
passages. These choices can coexist with content-derived structures for
finding passages. Table~\ref{tab:related-comparison} contrasts the objects
organised and searched by these systems.

\subsection{Agentic retrieval and source verification}

RAG draws on retrieved text during generation \citep{rag-lewis,gao-rag-survey-2023}.
Tool-using and adaptive retrieval workflows choose retrieval actions or
interleave search with reasoning \citep{react-yao,toolformer,ircot,adaptive-rag,modular-rag}.
EPRAG studies action selection across sources in enterprise RAG
\citep{kinalioglu-kbs-eprag-2026}; TEIE-RAG focuses on trustworthy engineering
information extraction \citep{shen-kbs-teierag-2026}. HMRS exposes catalogue,
shelf and source access through MCP \citep{mcp}; its contribution concerns
the managed state those actions use, not a new protocol or general agent
planner.

Neural rerankers score candidates to improve ordering
\citep{nogueira-monobert-2019,khattab-colbert-2020,sun-urara}. Self-RAG and
CRAG assess evidence or correct retrieval during generation
\citep{selfrag,crag}. HMRS distinguishes ordering documents from deciding
which passages enter an evidence pack; its optional gate scores source
text and returns admitted passages with references. A reference identifies
where to check the text, not whether it supports a generated claim
\citep{liu-verifiability,gao-citations}. Likewise, an interrupted search
must be distinguished from a completed search that found no support. Our
answer-quality comparison is separate from those retrieval checks, with
its judging protocol and limitations reported alongside the results
\citep{zheng-llmjudge,panickssery-selfpref,wang-llmnotfair,es-ragas-2024,saad-falcon-ares-2024}.

\subsection{Position of this work}
\label{sec:gap}

HMRS combines content-informed filing with collection and document
selection over the same maintained hierarchy. This gives the agent a way
from task scope to source passages and gives operators a way to check
whether a document remains visible after its managed state changes. The
contribution is this connection between organisation and retrieval, not the
invention of hierarchical search, scoped indexing or passage scoring.
